% Neue lokale Konstruktion. Die alten 17 Beweistabellen bleiben unverändert.
% Ausschließlich lokale Kurzzeichen des endlichen Zertifikats.
\providecommand{\FTList}[1]{\mathsf L(#1)}
\providecommand{\FTDistinct}[1]{\mathsf D(#1)}
\providecommand{\FTNum}[1]{\overline{#1}}
\providecommand{\FTCases}[1]{\FormulaRefAuto{FranklTripleFiniteCases}[thm]{#1}}
\providecommand{\FTProjection}[2]{\pi_{#1}(#2)}
\providecommand{\FTData}{\FTDistinct{a,b,c,d,e,f,g}}
% Reiner Zeilenumbruch eines verschachtelten Quellenverweises.
\providecommand{\FTSource}[2]{%
  \begin{gathered}[t]#1\bigl(\\[-2pt]#2\bigr)\end{gathered}}

\chapter{Eine enthaltene Dreiermenge ohne halbhäufiges Element}
\hypertarget{frankl.triple.construction}{}

Die Konstruktion hat 19 verschiedene Mengenglieder. Jedes der drei
Elemente der enthaltenen Menge \(\{a,b,c\}\) liegt in neun Gliedern und
fehlt in zehn. Wir führen den Nachweis auf endliche, ausdrücklich
angegebene Listen zurück. Die kurzen Regelschemata im ersten Abschnitt
begründen sowohl das Auswerten der Listen als auch den Größenvergleich;
eine Computerrechnung wird dabei nicht als Schlussregel verwendet.

\section{Endliche Listen als ausgeschriebene Mengenterme}

Die Indizes in diesem Abschnitt sind \emph{metasprachliche}, fest
vorgegebene endliche Laufbereiche. Eine Formel mit \(\bigwedge\) steht
für die angegebene endliche, rechtsgeklammerte Konjunktion; bei
\(\bigvee\) wird jeweils die letzte Alternative rechts an die bisherige
Disjunktion angefügt, beginnend mit \(D_0=\bot\):
\(D_{k+1}=D_k\lor P_{k+1}\). Damit bleibt die führende Falschheit
ausdrücklich erhalten. Beide Zeichen stehen für endliche Formeln,
nicht für einen zusätzlichen Quantor. Eine indizierte Tabellenzeile
bezeichnet die einzeln bezeichneten Instanzen derselben Herleitung.
Ihre Zusammenfassung erfolgt durch die in Band~01 eingeführte
\(\rAIRepeat{n}{\rho}{\sigma_1\mid\cdots\mid\sigma_n}\).
Insbesondere enthält jede Quellenangabe die Substitution und die
verwendeten vorherigen Zeilen.

\FormulaDefDeltaK[Endliche Listen und Verschiedenheitsbedingung]{%
  \begin{aligned}[t]
    \FTList{}&\coloneqq\varnothing,&
    \FTList{t_1,\ldots,t_n,t}&\coloneqq
      \FTList{t_1,\ldots,t_n}\cup\{t\},\\[-2pt]
    \FTDistinct{t_1,\ldots,t_n}&\coloneqq
      \bigwedge_{1\leq i<j\leq n}t_i\neq t_j,&
    \FTNum0&\coloneqq0,\quad
    \FTNum{n+1}\coloneqq\FTNum n+1.
  \end{aligned}
}{FranklTripleListNotation}{%
  \DeltaRow{Mengen}{t\dsep t_1\dsep\cdots\dsep t_n}
  \DeltaRow{Konvention}{\text{Leere Konjunktion: }\top;
     \quad\text{leere Disjunktion: }\bot}
}
Die Schreibweise \(\FTList{t_i:i\in I}\) ordnet einen ausdrücklich
angegebenen endlichen Indexbereich \(I\) aufsteigend. Wiederholungen
werden nicht vorausgesetzt; wo die Länge zählen soll, wird die
Verschiedenheit gesondert bewiesen. \(\pi_\alpha(r)\) bezeichnet die
Projektion des mit \(\alpha\) bezeichneten Konjunkts in Zeile \(r\).
Steht dieser Index an Position \(p\) der angegebenen geordneten
Indexliste der Länge \(N\), bedeutet dies für \(p<N\) genau
\(\land E_1\) nach \(p-1\) Anwendungen von \(\land E_2\), für
\(p=N\) nur diese \(N-1\) Anwendungen. Mehrfachindizes sind
lexikographisch geordnet; die Buchstabenfolge ist stets
\((a,b,c,d,e,f,g)\). Damit bezeichnet etwa \(\pi_{i,z}\) den
Eintrag der Zeile \(i\) und Spalte \(z\), auch wenn eine Teilliste
Lücken in ihren ursprünglichen Indizes hat.

\FormulaThmDeltaK[Die benötigten Zahlterme sind natürliche Zahlen]{%
  \bigwedge_{k=0}^{19}\FTNum k\in\mathbb N
}{FranklTripleNumeralsNatural}{\DeltaRow{Mengen}{\FTNum0\dsep\cdots\dsep\FTNum{19}}}
Die indizierte Zeile \(2(k)\) bezeichnet genau die 19 aufeinander
folgenden Instanzen für \(k=0,\ldots,18\); ihre Quellenrekursion
verweist ausschließlich auf die vorherige Instanz.
\begin{tabproofwide}
  \proofstepwidestar[]{\FTNum0\in\mathbb N}{\FormulaRefAuto{PeanoZero}[ax]}
  \proofstepwidestar[]{\FTNum{k+1}\in\mathbb N}{%
    \FormulaRefAuto{FranklTripleListNotation}[def]
      {\FormulaRefAuto{PeanoAddOneClosure}[thm]
        {1\ (k=0);\ 2(k-1)\ (1\leq k\leq18)}}}
  \proofstepwidestar[]{\bigwedge_{k=0}^{19}\FTNum k\in\mathbb N}{%
    \rAIStar{1,2(0),2(1),2(2),2(3),2(4),2(5),2(6),2(7),2(8),
      2(9),2(10),2(11),2(12),2(13),2(14),2(15),2(16),2(17),2(18)}}
\end{tabproofwide}

\FormulaThmDeltaK[Listen mit höchstens drei benannten Einträgen]{%
  \FTList{a}=\{a\}\land\FTList{a,b}=\{a,b\}\land
    \FTList{a,b,c}=\{a,b,c\}
}{FranklTripleSmallLists}{\DeltaRow{Mengen}{a\dsep b\dsep c}}
\begin{tabproofwide}
  \proofstepwidestar[]{\FTList{a}=\varnothing\cup\{a\}}{%
    \FormulaRefAuto{FranklTripleListNotation}[def]}
  \proofstepwidestar[]{\varnothing\cup\{a\}=\{a\}}{%
    \FormulaRefAuto{\varnothing\cup A=A}[thm]}
  \proofstepwidestar[]{\FTList{a}=\{a\}}{%
    \FormulaRefAuto{a=b\dsep b=c\vdash a=c}[thm]{1,2}}
  \proofstepwidestar[]{\FTList{a,b}=\FTList{a}\cup\{b\}}{%
    \FormulaRefAuto{FranklTripleListNotation}[def]}
  \proofstepwidestar[]{\FTList{a,b}=\{a\}\cup\{b\}}{\rIE{3,4}}
  \proofstepwidestar[]{\{a,b\}=\{a\}\cup\{b\}}{%
    \FormulaRefAuto{\{a,b\}=\{a\}\cup\{b\}}[thm]}
  \proofstepwidestar[]{\FTList{a,b}=\{a,b\}}{%
    \rIE{\FormulaRefAuto{a=b\vdash b=a}[thm]{6},5}}
  \proofstepwidestar[]{\FTList{a,b,c}=\FTList{a,b}\cup\{c\}}{%
    \FormulaRefAuto{FranklTripleListNotation}[def]}
  \proofstepwidestar[]{\FTList{a,b,c}=\{a,b\}\cup\{c\}}{\rIE{7,8}}
  \proofstepwidestar[]{\{a,b,c\}=\{a,b\}\cup\{c\}}{%
    \FormulaRefAuto{TripleSetDef}[def]}
  \proofstepwidestar[]{\FTList{a,b,c}=\{a,b,c\}}{%
    \rIE{\FormulaRefAuto{a=b\vdash b=a}[thm]{10},9}}
  \proofstepwidestar[]{\FTList{a}=\{a\}\land\FTList{a,b}=\{a,b\}\land
    \FTList{a,b,c}=\{a,b,c\}}{\rAIStar{3,7,11}}
\end{tabproofwide}

\FormulaThmDeltaK[Von zwei Herleitungen zum Bikonditional]{%
  [P]\vdots Q\dsep[Q]\vdots P\vdash P\leftrightarrow Q
}{FranklTripleEquivalenceRule}{\DeltaRow{Formeln}{P\dsep Q}}
\begin{tabproofwide}
  \proofstepwidestar[1]{P}{\rA}
  \proofstepwidestar[1,\Gamma]{Q}{%
    \begin{gathered}[t]\text{Erste angegebene}\\[-2pt]\text{Herleitung}\end{gathered}}
  \proofstepwidestar[\Gamma]{P\to Q}{\rRI{1,2}}
  \proofstepwidestar[4]{Q}{\rA}
  \proofstepwidestar[4,\Lambda]{P}{%
    \begin{gathered}[t]\text{Zweite angegebene}\\[-2pt]\text{Herleitung}\end{gathered}}
  \proofstepwidestar[\Lambda]{Q\to P}{\rRI{4,5}}
  \proofstepwidestar[\Gamma,\Lambda]{P\leftrightarrow Q}{\rLRI{3,6}}
\end{tabproofwide}
Ein Verweis auf \(P\eqvdash Q\) stellt genau diese beiden
Herleitungen bereit. Die vorliegende Regel macht daraus ausdrücklich
die objektsprachliche Formel \(P\leftrightarrow Q\).

\FormulaThmDeltaK[Mitgliedschaft in einer endlichen Liste]{%
  x\in\FTList{t_1,\ldots,t_n}
    \leftrightarrow\bigvee_{i=1}^{n}x=t_i
}{FranklTripleListMembership}{%
  \DeltaRow{Mengen}{x\dsep t_1\dsep\cdots\dsep t_n}
}
Der Nachweis ist eine Induktion über die Länge der \emph{geschriebenen}
Liste. Die geschlossene Anfangsherleitung lautet:
\begin{tabproofwide}
  \proofstepwidestar[1]{x\in\varnothing}{\rA}
  \proofstepwidestar[]{x\notin\varnothing}{%
    \FormulaRefAuto{x\not\in\varnothing}[thm]}
  \proofstepwidestar[1]{\bot}{\rBI{1,2}}
  \proofstepwidestar[]{x\in\varnothing\to\bot}{\rRI{1,3}}
  \proofstepwidestar[5]{\bot}{\rA}
  \proofstepwidestar[5]{x\in\varnothing}{%
    \FormulaRefAuto{\bot\vdash P}[thm]{5}}
  \proofstepwidestar[]{\bot\to x\in\varnothing}{\rRI{5,6}}
  \proofstepwidestar[]{x\in\FTList{}\leftrightarrow\bot}{%
    \FormulaRefAuto{FranklTripleListNotation}[def]{\rLRI{4,7}}}
\end{tabproofwide}
Der folgende Erweiterungsschritt beweist jede weitere Instanz. Dabei ist
\(L=\FTList{t_1,\ldots,t_n}\) und \(Q=\bigvee_{i=1}^n x=t_i\).
\begin{tabproofwide}
  \proofstepwidestar[1]{x\in L\leftrightarrow Q}{\rA}
  \proofstepwidestar[]{x\in L\cup\{t\}\leftrightarrow
    (x\in L\lor x\in\{t\})}{%
    \FormulaRefAuto{FranklTripleEquivalenceRule}[thm]
      {\FormulaRefAuto{z\in A\cup B\eqvdash z\in A\lor z\in B}[thm]}}
  \proofstepwidestar[]{x\in\{t\}\leftrightarrow x=t}{%
    \FormulaRefAuto{FranklTripleEquivalenceRule}[thm]
      {\FormulaRefAuto{x\in\{a\}\eqvdash x=a}[thm]}}
  \proofstepwidestar[1]{x\in L\cup\{t\}\leftrightarrow(Q\lor x=t)}{%
    \rLRS{1,\rLRS{3,2}}}
  \proofstepwidestar[1]{x\in\FTList{t_1,\ldots,t_n,t}
    \leftrightarrow\bigl(Q\lor x=t\bigr)}{%
    \FormulaRefAuto{FranklTripleListNotation}[def]{4}}
\end{tabproofwide}
Die hier verwendete Verknüpfung der letzten Alternative mit der
bisherigen Disjunktion legt deren Klammerung fest; die unten verwendete
Fallregel folgt genau dieser Klammerung. Die offene Annahme des
Erweiterungsschritts wird jeweils durch die bereits geschlossene
Vorgängerherleitung ersetzt.

\FormulaThmDeltaK[Einführung eines aufgeführten Elements]{%
  x=t_i\vdash x\in\FTList{t_1,\ldots,t_n}
}{FranklTripleListIntroduction}{%
  \DeltaRow{Mengen}{x\dsep t_1\dsep\cdots\dsep t_n}
  \DeltaRow{Index}{1\leq i\leq n}
}
Der Quellenkurzterm \(\lambda_{i,n}(r)\) enthält ausschließlich
Disjunktionseinführungen:
\[
  \lambda_{i,i}(r)=\rOIb{r},\qquad
  \lambda_{i,k+1}(r)=\rOIa{\lambda_{i,k}(r)}\quad(i\leq k<n).
\]
Die erste Anwendung führt die rechte Alternative von
\(D_{i-1}\lor x=t_i\) ein; die übrigen fügen die jeweils nächste
Alternative rechts an. So sind auch der erste und letzte Listenplatz
ohne Sonderregel erfasst.
\begin{tabproofwide}
  \proofstepwidestar[1]{x=t_i}{\rA}
  \proofstepwidestar[1]{\bigvee_{j=1}^{n}x=t_j}{%
    \lambda_{i,n}(1)}
  \proofstepwidestar[]{x\in\FTList{t_1,\ldots,t_n}
      \leftrightarrow\bigvee_{j=1}^{n}x=t_j}{%
    \FormulaRefAuto{FranklTripleListMembership}[thm]}
  \proofstepwidestar[1]{x\in\FTList{t_1,\ldots,t_n}}{%
    \FormulaRefAuto{P\leftrightarrow Q,Q\vdash P}[thm]{3,2}}
\end{tabproofwide}

\FormulaThmDeltaK[Endliche Fallregel]{%
  \begin{aligned}[t]
  &x\in\FTList{t_1,\ldots,t_n}\dsep
    [x=t_1]\vdots R\dsep\cdots\dsep[x=t_n]\vdots R\vdash R
  \end{aligned}
}{FranklTripleFiniteCases}{%
  \DeltaRow{Mengen}{x\dsep t_1\dsep\cdots\dsep t_n}
  \DeltaRow{Formel}{R}
}
Nach \FormulaRefAuto{FranklTripleListMembership}[thm] genügt die
gewöhnliche Disjunktionsbeseitigung. Für null Alternativen folgt
\(R\) aus \(\bot\) durch \FormulaRefAuto{\bot\vdash P}[thm].
Der dafür verwendete Regelterm ist vollständig rekursiv festgelegt:
\[
\begin{aligned}
 \mathcal E_0(r)&=\FormulaRefAuto{\bot\vdash P}[thm]{r},\\
 \mathcal E_{n+1}(r;\rho_1,\ldots,\rho_{n+1})
   &=\rOE{r,\alpha,\mathcal E_n(\alpha;\rho_1,\ldots,\rho_n),
       \beta,\rho_{n+1}(\beta)}.
\end{aligned}
\]
Dabei ist \(\alpha\) die neue Annahme der bisherigen Disjunktion und
\(\beta\) die Annahme der letzten Alternative; \(\rho_i\) ist die
jeweils angegebene Fallherleitung. Beim Aufruf mit einer
Listenmitgliedschaft in Zeile \(r\) wird als erstes Argument
\[
 \FormulaRefAuto{P\leftrightarrow Q,P\vdash Q}[thm]
   {\FormulaRefAuto{FranklTripleListMembership}[thm],r}
\]
eingesetzt. Der folgende Schritt fügt genau eine
Alternative hinzu und entlädt ausschließlich deren Fallannahmen.
\begin{tabproofwide}
  \proofstepwidestar[1]{Q\lor x=t}{\rA}
  \proofstepwidestar[2]{Q}{\rA}
  \proofstepwidestar[2,\Gamma]{R}{%
    \mathcal E_n(2;\rho_1,\ldots,\rho_n)}
  \proofstepwidestar[4]{x=t}{\rA}
  \proofstepwidestar[4,\Lambda]{R}{%
    \rho_{n+1}(4)}
  \proofstepwidestar[1,\Gamma,\Lambda]{R}{\rOE{1,2,3,4,5}}
\end{tabproofwide}
Hier sind \(\Gamma,\Lambda\) die übrigen offenen Annahmen der
eingesetzten Herleitungen. Die Quellenangabe einer Anwendung nennt
stets die Listenzeile und alle Fallherleitungen, gegebenenfalls durch
einen ausdrücklich festgelegten endlichen Indexbereich.

\FormulaThmDeltaK[Vereinigung von Listen]{%
  \FTList{t_i:i\in I}\cup\FTList{t_i:i\in J}
    =\FTList{t_i:i\in I\cup J}
}{FranklTripleUnionNormalForm}{%
  \DeltaRow{Mengen}{t_i\dsep x}
  \DeltaRow{Indexbereiche}{I,J\text{ sind endliche Indexlisten ohne Wiederholung}}
}
Die Mengenzeichen an den Indexbereichen beschreiben lediglich,
welche geschriebenen Listeneinträge übernommen werden. Setze
\(L_I=\FTList{t_i:i\in I}\), entsprechend \(L_J,L_{I\cup J}\).
\begin{tabproofsplitwide}
  \proofpartwide{Von links nach rechts}
  \proofstepwidestar[1]{x\in L_I\cup L_J}{\rA}
  \proofstepwidestar[1]{x\in L_I\lor x\in L_J}{%
    \FormulaRefAuto{z\in A\cup B\eqvdash z\in A\lor z\in B}[thm]{1}}
  \proofstepwidestar[3]{x\in L_I}{\rA}
  \proofstepwidestar[4]{x=t_i}{\rA\quad(i\in I)}
  \proofstepwidestar[4]{x\in L_{I\cup J}}{%
    \FormulaRefAuto{FranklTripleListIntroduction}[thm]{4}}
  \proofstepwidestar[3]{x\in L_{I\cup J}}{%
    \FTCases{3;\ [4]\vdots5\ (i\in I)}}
  \proofstepwidestar[7]{x\in L_J}{\rA}
  \proofstepwidestar[8]{x=t_j}{\rA\quad(j\in J)}
  \proofstepwidestar[8]{x\in L_{I\cup J}}{%
    \FormulaRefAuto{FranklTripleListIntroduction}[thm]{8}}
  \proofstepwidestar[7]{x\in L_{I\cup J}}{%
    \FTCases{7;\ [8]\vdots9\ (j\in J)}}
  \proofstepwidestar[1]{x\in L_{I\cup J}}{\rOE{2,3,6,7,10}}
  \proofstepwidestar[]{L_I\cup L_J\subseteq L_{I\cup J}}{%
    \FormulaRefAuto{SubsetIntroductionRule}[thm]{[1]\vdots11}}
  \closeproofpartwide
  \proofpartwide{Von rechts nach links}
  \setcounter{proofstepnr}{12}
  \proofstepwidestar[13]{x\in L_{I\cup J}}{\rA}
  \proofstepwidestar[14]{x=t_k}{\rA\quad(k\in I\cup J)}
  \proofstepwidestar[14]{x\in L_I\cup L_J}{%
    \begin{gathered}[t]
      k\in I:\\[-2pt]
      \FTSource{\FormulaRefAuto{z\in A\vdash z\in A\cup B}[thm]}
        {\FormulaRefAuto{FranklTripleListIntroduction}[thm]{14}}
          ,\\[2pt]
      k\in J\setminus I:\\[-2pt]
      \FTSource{\FormulaRefAuto{z\in B\vdash z\in A\cup B}[thm]}
        {\FormulaRefAuto{FranklTripleListIntroduction}[thm]{14}}
    \end{gathered}}
  \proofstepwidestar[13]{x\in L_I\cup L_J}{%
    \FTCases{13;\ [14]\vdots15\ (k\in I\cup J)}}
  \proofstepwidestar[]{L_{I\cup J}\subseteq L_I\cup L_J}{%
    \FormulaRefAuto{SubsetIntroductionRule}[thm]{[13]\vdots16}}
  \proofstepwidestar[]{L_I\cup L_J=L_{I\cup J}}{%
    \FormulaRefAuto{A\subseteq B,B\subseteq A\vdash A=B}[thm]{12,17}}
  \closeproofpartwide
\end{tabproofsplitwide}

\FormulaThmDeltaK[Mitgliedschafts- und Nichtmitgliedschaftszeugnisse]{%
  \FTDistinct{t_1,\ldots,t_n}\vdash
  \begin{cases}
    t_j\in\FTList{t_i:i\in I},&j\in I,\\[-2pt]
    t_j\notin\FTList{t_i:i\in I},&j\notin I
  \end{cases}
}{FranklTripleIncidence}{%
  \DeltaRow{Mengen}{t_1\dsep\cdots\dsep t_n}
  \DeltaRow{Indexbereiche}{I\subseteq\{1,\ldots,n\},\quad1\leq j\leq n}
}
Der erste Fall folgt aus \(t_j=t_j\) durch
\FormulaRefAuto{FranklTripleListIntroduction}[thm]. Im zweiten Fall
liefert jede Verschiedenheitsprojektion, nötigenfalls nach Symmetrie
der angenommenen Gleichheit, den folgenden Widerspruch.
\begin{tabproofwide}
  \proofstepwidestar[1]{\FTDistinct{t_1,\ldots,t_n}}{\rA}
  \proofstepwidestar[2]{t_j\in\FTList{t_i:i\in I}}{\rA}
  \proofstepwidestar[3]{t_j=t_i}{\rA\quad(i\in I)}
  \proofstepwidestar[1,3]{\bot}{%
    \rBI{3,\FTProjection{(j,i)}{1}}\ (j<i);\quad
    \rBI{\FormulaRefAuto{a=b\vdash b=a}[thm]{3},
      \FTProjection{(i,j)}{1}}\ (i<j)}
  \proofstepwidestar[1,2]{\bot}{\FTCases{2;\ [3]\vdots4\ (i\in I)}}
  \proofstepwidestar[1]{t_j\notin\FTList{t_i:i\in I}}{\rCI{2,5}}
\end{tabproofwide}

\FormulaThmDeltaK[Eine verschiedene Liste zählt ihre Länge]{%
  \FTDistinct{t_1,\ldots,t_n}\vdash
  \NatSeg{\FTNum n}\EqCard\FTList{t_1,\ldots,t_n}
}{FranklTripleListCount}{%
  \DeltaRow{Mengen}{t_1\dsep\cdots\dsep t_n}
  \DeltaRow{Listenlänge}{0\leq n\leq19}
}
Wieder handelt es sich um eine Induktion über die geschriebene
Listenlänge. Ihre geschlossene Anfangsherleitung ist:
\begin{tabproofwide}
  \proofstepwidestar[]{\NatSeg{\FTNum0}=\varnothing}{%
    \FormulaRefAuto{FranklTripleListNotation}[def]
      {\FormulaRefAuto{PeanoNatSegZero}[thm]}}
  \proofstepwidestar[]{\varnothing\EqCard\varnothing}{%
    \FormulaRefAuto{EqCardReflexive}[thm]}
  \proofstepwidestar[]{\NatSeg{\FTNum0}\EqCard\FTList{}}{%
    \FormulaRefAuto{FranklTripleListNotation}[def]
      {\rIE{\FormulaRefAuto{a=b\vdash b=a}[thm]{1},2}}}
\end{tabproofwide}
Im Erweiterungsschritt seien \(L=\FTList{t_1,\ldots,t_n}\)
und \(t=t_{n+1}\), wobei \(0\leq n<19\).
Für \(n=0,1\) wird in Zeile~2 der geschlossene Satz
\FormulaRefAuto{\top}[thm] benutzt; es gibt keine Wiederholung mit
null Instanzen. Die ausdrücklich notierte Konjunktionseinführung mit
Zeile~1 und anschließende Projektion erhält deren Abhängigkeit.
Für \(n=0\) wird Zeile~4 analog aus
\FormulaRefAuto{x\not\in\varnothing}[thm] gewonnen.
Schreibe \(\rho_0\) für die geschlossene dritte Zeile der
Anfangsherleitung. Der Quellenkurzterm
\(\mathcal C_0(r)=\rAEa{\rAI{\rho_0,r}}\) nimmt eine
Verschiedenheitsprämisse auf. Rekursiv bezeichnet
\(\mathcal C_{n+1}(r)\) die folgende neunzeilig aufgeführte
Herleitung mit \(r\) an Stelle der Annahme~1 und
\(\mathcal C_n(2)\) in Zeile~3. Diese Definition verweist nur auf
die schon konstruierte kürzere Ableitung.
Die beiden längeren Quellen des Erweiterungsschritts schreiben wir
vorab aus; dies ändert nur ihren Platz auf der Seite:
\[
\begin{aligned}
 \beta_n(r)&=
 \begin{cases}
   \rAIRepeat{n(n-1)/2}{\operatorname{id}}
     {\pi_{(i,j)}(r)\ (1\leq i<j\leq n)},&n\geq2,\\
   \rAEa{\rAI{\FormulaRefAuto{\top}[thm],r}},&n=0,1;
 \end{cases}\\[3pt]
 \gamma_n(r)&=
 \begin{cases}
   \FormulaRefAuto{FranklTripleIncidence}[thm]{r},&n\geq1,\\
   \rAEa{\rAI{\FormulaRefAuto{x\not\in\varnothing}[thm],r}},&n=0.
 \end{cases}
\end{aligned}
\]
\begin{tabproofwide}
  \proofstepwidestar[1]{\FTDistinct{t_1,\ldots,t_n,t}}{\rA}
  \proofstepwidestar[1]{\FTDistinct{t_1,\ldots,t_n}}{\beta_n(1)}
  \proofstepwidestar[1]{\NatSeg{\FTNum n}\EqCard L}{%
    \mathcal C_n(2)}
  \proofstepwidestar[1]{t\notin L}{\gamma_n(1)}
  \proofstepwidestar[]{\FTNum n\in\mathbb N}{%
    \pi_n(\FormulaRefAuto{FranklTripleNumeralsNatural}[thm])}
  \proofstepwidestar[]{\FTNum n\notin\NatSeg{\FTNum n}}{%
    \FormulaRefAuto{PeanoNatSegEndpointNotIn}[thm]
      {5}}
  \proofstepwidestar[1]{\NatSeg{\FTNum n}\cup\{\FTNum n\}
       \EqCard L\cup\{t\}}{%
    \FormulaRefAuto{EqCardFreshAdjunction}[thm]{3,6,4}}
  \proofstepwidestar[]{\NatSeg{\FTNum{n+1}}
       =\NatSeg{\FTNum n}\cup\{\FTNum n\}}{%
    \FormulaRefAuto{PeanoNatSegAddOne}[thm]{5}}
  \proofstepwidestar[1]{\NatSeg{\FTNum{n+1}}
       \EqCard\FTList{t_1,\ldots,t_n,t}}{%
    \FormulaRefAuto{FranklTripleListNotation}[def]{\rIE{
      \FormulaRefAuto{a=b\vdash b=a}[thm]{8},7}}}
\end{tabproofwide}

\FormulaThmDeltaK[Endliches Vereinigungszertifikat]{%
  \bigwedge_{1\leq i,j\leq n}(X_i\cup X_j=X_{k(i,j)})
    \vdash\UnionClosedFam(\FTList{X_1,\ldots,X_n})
}{FranklTripleClosureCertificate}{%
  \DeltaRow{Mengen}{X_1\dsep\cdots\dsep X_n\dsep X\dsep Y}
  \DeltaRow{Indexvorschrift}{1\leq k(i,j)\leq n\text{ für jedes }i,j}
}
\begin{tabproofwide}
  \proofstepwidestar[1]{\bigwedge_{i,j}(X_i\cup X_j=X_{k(i,j)})}{\rA}
  \proofstepwidestar[2]{X\in\FTList{X_1,\ldots,X_n}}{\rA}
  \proofstepwidestar[3]{Y\in\FTList{X_1,\ldots,X_n}}{\rA}
  \proofstepwidestar[4]{X=X_i}{\rA\quad(1\leq i\leq n)}
  \proofstepwidestar[5]{Y=X_j}{\rA\quad(1\leq j\leq n)}
  \proofstepwidestar[1,4,5]{X\cup Y=X_{k(i,j)}}{%
    \rIE{\FormulaRefAuto{a=b\vdash b=a}[thm]{4},
          \FormulaRefAuto{a=b\vdash b=a}[thm]{5},\pi_{(i,j)}(1)}}
  \proofstepwidestar[1,4,5]{X\cup Y\in\FTList{X_1,\ldots,X_n}}{%
    \FormulaRefAuto{FranklTripleListIntroduction}[thm]{6}}
  \proofstepwidestar[1,3,4]{X\cup Y\in\FTList{X_1,\ldots,X_n}}{%
    \FTCases{3;\ [5]\vdots7\ (1\leq j\leq n)}}
  \proofstepwidestar[1,2,3]{X\cup Y\in\FTList{X_1,\ldots,X_n}}{%
    \FTCases{2;\ [4]\vdots8\ (1\leq i\leq n)}}
  \proofstepwidestar[1]{\UnionClosedFam(\FTList{X_1,\ldots,X_n})}{%
    \FormulaRefAuto{\UnionClosedFam(M)}[def]
      {\rUI{\rUI{\rRI{2,\rRI{3,9}}}}}}
\end{tabproofwide}

\FormulaThmDeltaK[Aussonderung aus einer endlichen Liste]{%
  \begin{aligned}[t]
  &\bigwedge_{i\in I}P(X_i)\land
      \bigwedge_{i\notin I}\neg P(X_i)\vdash{}\\[-2pt]
  &\{X\in\FTList{X_1,\ldots,X_n}\mid P(X)\}
      =\FTList{X_i:i\in I}
  \end{aligned}
}{FranklTripleFilterCertificate}{%
  \DeltaRow{Mengen}{X_1\dsep\cdots\dsep X_n\dsep X}
  \DeltaRow{Prädikat}{P}
  \DeltaRow{Indexbereich}{I\subseteq\{1,\ldots,n\}}
}
Setze \(L=\FTList{X_1,\ldots,X_n}\), \(L_I=\FTList{X_i:i\in I}\)
und \(S=\{X\in L\mid P(X)\}\).
\begin{tabproofsplitwide}
  \proofpartwide{Die Aussonderung enthält nur die aufgeführten Glieder}
  \proofstepwidestar[1]{\bigwedge_{i\in I}P(X_i)\land
      \bigwedge_{i\notin I}\neg P(X_i)}{\rA}
  \proofstepwidestar[2]{X\in S}{\rA}
  \proofstepwidestar[2]{X\in L\land P(X)}{%
    \FormulaRefAuto{x\in\{u\in A\mid P(u)\}\eqvdash x\in A\land P(x)}[thm]{2}}
  \proofstepwidestar[4]{X=X_j}{\rA\quad(1\leq j\leq n)}
  \proofstepwidestar[4]{X\in L_I}{%
    \FormulaRefAuto{FranklTripleListIntroduction}[thm]{4}
       \quad(j\in I)}
  \proofstepwidestar[2,4]{P(X_j)}{\rIE{4,\rAEb{3}}\quad(j\notin I)}
  \proofstepwidestar[1,2,4]{\bot}{%
    \rBI{6,\pi_j(\rAEb{1})}\quad(j\notin I)}
  \proofstepwidestar[1,2,4]{X\in L_I}{%
    \FormulaRefAuto{\bot\vdash P}[thm]{7}\quad(j\notin I)}
  \proofstepwidestar[1,2]{X\in L_I}{%
    \FTCases{\rAEa{3};[4]\vdots5\ (j\in I);[4]\vdots8\ (j\notin I)}}
  \proofstepwidestar[1]{S\subseteq L_I}{%
    \FormulaRefAuto{SubsetIntroductionRule}[thm]{[2]\vdots9}}
  \closeproofpartwide
  \proofpartwide{Jedes aufgeführte Glied liegt in der Aussonderung}
  \setcounter{proofstepnr}{10}
  \proofstepwidestar[11]{X\in L_I}{\rA}
  \proofstepwidestar[12]{X=X_i}{\rA\quad(i\in I)}
  \proofstepwidestar[12]{X\in L}{%
    \FormulaRefAuto{FranklTripleListIntroduction}[thm]{12}}
  \proofstepwidestar[1,12]{P(X)}{%
    \rIE{\FormulaRefAuto{a=b\vdash b=a}[thm]{12},\pi_i(\rAEa{1})}}
  \proofstepwidestar[1,12]{X\in S}{%
    \FormulaRefAuto{x\in A\dsep P(x)\vdash x\in\{x\in A\mid P(x)\}}[thm]{13,14}}
  \proofstepwidestar[1,11]{X\in S}{\FTCases{11;[12]\vdots15\ (i\in I)}}
  \proofstepwidestar[1]{L_I\subseteq S}{%
    \FormulaRefAuto{SubsetIntroductionRule}[thm]{[11]\vdots16}}
  \proofstepwidestar[1]{S=L_I}{%
    \FormulaRefAuto{A\subseteq B,B\subseteq A\vdash A=B}[thm]{10,17}}
  \closeproofpartwide
\end{tabproofsplitwide}

\clearpage
\section{Die neunzehn ausdrücklich angegebenen Mengenglieder}

\FormulaDefDeltaK[Die Familie des Beispiels]{%
  M_0\coloneqq\FTList{X_1,\ldots,X_{19}}
}{FranklTripleExampleDefinition}{%
  \DeltaRow{Mengen}{M_0\dsep a\dsep b\dsep c\dsep d\dsep e\dsep f\dsep g}
  \DeltaRow{Konvention}{X_i=\FTList{\text{die mit 1 markierten Einträge der Zeile }i}}
}
Die folgende Tabelle ist Teil dieser Definition; jede Zeile ist ein
konkreter, durch \FormulaRefAuto{FranklTripleListNotation}[def]
expandierbarer Mengenterm. Die Blockbezeichnungen fassen nur Zeilen
zusammen. Es gilt \(A=\FTList{a,b,c}=\{a,b,c\}\),
\(B=\FTList{d,e,f,g}\), \(D=\FTList{d,e,g}\),
\(E=\FTList{d,f,g}\), \(F=\FTList{e,f,g}\).
\begin{center}\small
\begin{tabular}{>{$}c<{$}>{$}c<{$}>{$}l<{$}ccccccc}
\toprule
\text{Block}&i&X_i&a&b&c&d&e&f&g\\
\midrule
H_0&1&\varnothing&0&0&0&0&0&0&0\\
&2&A&1&1&1&0&0&0&0\\
\midrule
H_a&3&D&0&0&0&1&1&0&1\\
&4&D\cup\{a\}&1&0&0&1&1&0&1\\
&5&D\cup A&1&1&1&1&1&0&1\\
\midrule
H_b&6&E&0&0&0&1&0&1&1\\
&7&E\cup\{b\}&0&1&0&1&0&1&1\\
&8&E\cup A&1&1&1&1&0&1&1\\
\midrule
H_c&9&F&0&0&0&0&1&1&1\\
&10&F\cup\{c\}&0&0&1&0&1&1&1\\
&11&F\cup A&1&1&1&0&1&1&1\\
\midrule
H_*&12&B&0&0&0&1&1&1&1\\
&13&B\cup\{a\}&1&0&0&1&1&1&1\\
&14&B\cup\{b\}&0&1&0&1&1&1&1\\
&15&B\cup\{c\}&0&0&1&1&1&1&1\\
&16&B\cup\{a,b\}&1&1&0&1&1&1&1\\
&17&B\cup\{a,c\}&1&0&1&1&1&1&1\\
&18&B\cup\{b,c\}&0&1&1&1&1&1&1\\
&19&B\cup A&1&1&1&1&1&1&1\\
\bottomrule
\end{tabular}
\end{center}
Die ausgeschriebenen Namen in der dritten Spalte stimmen nach
\FormulaRefAuto{FranklTripleUnionNormalForm}[thm] mit den in den
sieben rechten Spalten definierten Listentermen überein.

\FormulaThmDeltaK[Die Einträge der Inzidenztabelle gelten]{%
  \FTData\vdash
  \bigwedge_{(i,z)} I_{iz}
}{FranklTripleExampleIncidence}{%
  \DeltaRow{Mengen}{a\dsep b\dsep c\dsep d\dsep e\dsep f\dsep g\dsep X_i}
  \DeltaRow{Formel}{I_{iz}\text{ ist }z\in X_i\text{ bei 1 und }z\notin X_i\text{ bei 0}}
  \DeltaRow{Indexbereich}{1\leq i\leq19,\quad z\in(a,b,c,d,e,f,g)}
}
Die Substitution \(s_{iz}\) setzt die sieben \(t\)-Variablen auf
\((a,b,c,d,e,f,g)\), den Index \(j\) auf die Spaltenposition von
\(z\) und \(I\) auf die mit 1 markierten Spalten der Zeile \(i\).
\begin{tabproofwide}
  \proofstepwidestar[1]{\FTData}{\rA}
  \proofstepwidestar[1]{\bigwedge_{(i,z)} I_{iz}}{%
    \rAIRepeat{133}{\FormulaRefAuto{FranklTripleIncidence}[thm]}{
      1;\ s_{iz}\ (1\leq i\leq19)}}
\end{tabproofwide}

\FormulaThmDeltaK[Die Mengenglieder sind paarweise verschieden]{%
  \FTData\vdash\FTDistinct{X_1,\ldots,X_{19}}
}{FranklTripleExampleDistinct}{%
  \DeltaRow{Mengen}{a\dsep b\dsep c\dsep d\dsep e\dsep f\dsep g\dsep X_i}
}
Für jedes \(i<j\) wähle als \(w(i,j)\) die erste Spalte, in der die
beiden geschriebenen Zeilen verschieden sind. Die Tabelle hat keine
zwei gleichen Zeilen: verschiedene Blöcke unterscheiden sich in
\(d,e,f,g\), innerhalb eines Blocks unterscheiden sich die Muster in
\(a,b,c\). Dies legt alle 171 Zeugen ohne Wahlannahme eindeutig fest.
Für die jeweilige Inzidenzrichtung ist der Widerspruchsterm
\(\omega_{ij}(r,s)\) vollständig festgelegt; hierbei steht \(r\)
für die Inzidenzen, \(s\) für \(X_i=X_j\) und \(w=w(i,j)\):
\[
\omega_{ij}(r,s)=
\begin{cases}
 \rBI{\rIE{s,\pi_{i,w}(r)},\pi_{j,w}(r)},&I_{i,w}\text{ ist }w\in X_i,\\[2pt]
 \rBI{\rIE{\FormulaRefAuto{a=b\vdash b=a}[thm]{s},\pi_{j,w}(r)},
       \pi_{i,w}(r)},&I_{j,w}\text{ ist }w\in X_j.
\end{cases}
\]
\begin{tabproofwide}
  \proofstepwidestar[1]{\FTData}{\rA}
  \proofstepwidestar[1]{\bigwedge_{(i,z)}I_{iz}}{%
    \FormulaRefAuto{FranklTripleExampleIncidence}[thm]{1}}
  \proofstepwidestar[3]{X_i=X_j}{\rA\quad(i<j)}
  \proofstepwidestar[1,3]{\bot}{\omega_{ij}(2,3)}
  \proofstepwidestar[1]{X_i\neq X_j}{\rCI{3,4}\quad(i<j)}
  \proofstepwidestar[1]{\FTDistinct{X_1,\ldots,X_{19}}}{%
    \rAIRepeat{171}{\neg I}{[3]\vdots4\ (1\leq i<j\leq19)}}
\end{tabproofwide}

\FormulaThmDeltaK[Das Vereinigungszertifikat der fünf Blöcke]{%
  \bigwedge_{1\leq i,j\leq19}(X_i\cup X_j=X_{k(i,j)})
}{FranklTripleExampleUnions}{%
  \DeltaRow{Mengen}{a\dsep b\dsep c\dsep d\dsep e\dsep f\dsep g\dsep X_i}
}
Die Zielvorschrift \(k(i,j)\) lässt sich ohne 361 einzelne Felder
angeben. Der Zielblock steht in der folgenden kleinen Tafel. Innerhalb
dieses Blocks werden die drei \(a,b,c\)-Spalten durch das logische
\emph{Oder} der beiden Eingabezeilen bestimmt.
\begin{center}\small
\begin{tabular}{>{$}c<{$}|*{5}{>{$}c<{$}}}
 &H_0&H_a&H_b&H_c&H_*\\\hline
H_0&H_0&H_a&H_b&H_c&H_*\\
H_a&H_a&H_a&H_*&H_*&H_*\\
H_b&H_b&H_*&H_b&H_*&H_*\\
H_c&H_c&H_*&H_*&H_c&H_*\\
H_*&H_*&H_*&H_*&H_*&H_*\\
\end{tabular}
\end{center}
Diese Vorschrift benennt immer genau eine vorhandene Zeile.
In \(H_0\) sind nur \(000\) und \(111\) zugelassen.
Die drei mittleren Blöcke erlauben der Reihe nach die Muster
\(000,100,111\), \(000,010,111\) und \(000,001,111\).
Jede dieser Listen ist unter Oder abgeschlossen; Oder mit \(000\)
oder \(111\) bleibt darin. Der Block \(H_*\) enthält alle acht Muster.
Die vier letzten Spalten ergeben genau die dargestellte Blocktafel:
die verschiedenen mittleren Muster \(1101,1011,0111\) ergeben
paarweise \(1111\), während \(0000\) neutral und \(1111\)
absorbierend ist. Damit ist für jedes benannte \(i,j\) die
Vereinigung ihrer beiden expliziten Indexlisten genau die Indexliste
von \(k(i,j)\). Der eigentliche Mengenschluss ist die bereits
bewiesene Listenregel:
\begin{tabproofwide}
  \proofstepwidestar[]{X_i\cup X_j=X_{k(i,j)}}{%
    \begin{gathered}[t]
      \FTSource{\FormulaRefAuto{FranklTripleExampleDefinition}[def]}
        {\FormulaRefAuto{FranklTripleUnionNormalForm}[thm]}\\[-2pt]
      (1\leq i,j\leq19)
    \end{gathered}}
  \proofstepwidestar[]{\bigwedge_{1\leq i,j\leq19}
      (X_i\cup X_j=X_{k(i,j)})}{%
    \rAIRepeat{361}{\operatorname{id}}{1(i,j)\ (1\leq i,j\leq19)}}
\end{tabproofwide}

\FormulaThmDeltaK[Endlichkeit, Abschluss und enthaltene Dreiermenge]{%
  \FTData\vdash\Finite{M_0}\land\UnionClosedFam(M_0)
       \land\{a,b,c\}\in M_0
}{FranklTripleExampleFamily}{%
  \DeltaRow{Mengen}{M_0\dsep a\dsep b\dsep c\dsep d\dsep e\dsep f\dsep g}
}
\begin{tabproofwide}
  \proofstepwidestar[1]{\FTData}{\rA}
  \proofstepwidestar[1]{\FTDistinct{X_1,\ldots,X_{19}}}{%
    \FormulaRefAuto{FranklTripleExampleDistinct}[thm]{1}}
  \proofstepwidestar[1]{\NatSeg{\FTNum{19}}\EqCard M_0}{%
    \FormulaRefAuto{FranklTripleExampleDefinition}[def]
      {\FormulaRefAuto{FranklTripleListCount}[thm]{2}}}
  \proofstepwidestar[]{\FTNum{19}\in\mathbb N}{%
    \pi_{19}(\FormulaRefAuto{FranklTripleNumeralsNatural}[thm])}
  \proofstepwidestar[1]{\Finite{M_0}}{%
    \FormulaRefAuto{FiniteNatSegEqCardIntro}[thm]{4,3}}
  \proofstepwidestar[]{\UnionClosedFam(M_0)}{%
    \FormulaRefAuto{FranklTripleExampleDefinition}[def]
      {\FormulaRefAuto{FranklTripleClosureCertificate}[thm]
      {\FormulaRefAuto{FranklTripleExampleUnions}[thm]}}}
  \proofstepwidestar[]{\{a,b,c\}=X_2}{%
    \FormulaRefAuto{FranklTripleExampleDefinition}[def]
      {\FormulaRefAuto{a=b\vdash b=a}[thm]
        {\rAEb{\rAEb{\FormulaRefAuto{FranklTripleSmallLists}[thm]}}}}}
  \proofstepwidestar[]{\{a,b,c\}\in M_0}{%
    \FormulaRefAuto{FranklTripleExampleDefinition}[def]
      {\FormulaRefAuto{FranklTripleListIntroduction}[thm]{7}}}
  \proofstepwidestar[1]{\begin{aligned}[t]
       &\Finite{M_0}\land\UnionClosedFam(M_0)\\[-2pt]
       &{}\land\{a,b,c\}\in M_0
     \end{aligned}}{\rAIStar{5,6,8}}
\end{tabproofwide}

\section{Neun enthaltende und zehn vermeidende Glieder}

Die Indexlisten der beiden Teilfamilien werden vollständig angegeben.
Ihre Längen sind die ausgeschriebenen neun beziehungsweise zehn
Positionen; die Verschiedenheit folgt aus
\FormulaRefAuto{FranklTripleExampleDistinct}[thm].
\begin{center}\small
\begin{tabular}{>{$}c<{$}>{$}l<{$}>{$}l<{$}}
\toprule
z&I_z^+&I_z^-\\\midrule
a&2,4,5,8,11,13,16,17,19&1,3,6,7,9,10,12,14,15,18\\
b&2,5,7,8,11,14,16,18,19&1,3,4,6,9,10,12,13,15,17\\
c&2,5,8,10,11,15,17,18,19&1,3,4,6,7,9,12,13,14,16\\
\bottomrule
\end{tabular}
\end{center}

Die Quellen für die Aussonderungsregel werden ausdrücklich
zusammengesetzt. Für eine Zeile \(r\), die sämtliche Inzidenzen
enthält, gelten die folgenden Kurzformen für Regelterme:
\[
\begin{aligned}
 \alpha_z^+(r)&=\rAIRepeat{9}{\operatorname{id}}
      {\pi_{i,z}(r)\ (i\in I_z^+)},\\
 \alpha_z^-(r)&=\rAIRepeat{10}{\operatorname{id}}
      {\pi_{i,z}(r)\ (i\in I_z^-)},\\
 \widehat\alpha_z^+(r)&=\rAIRepeat{9}{\FormulaRefAuto{rDN}[thm]}
      {\pi_{i,z}(r)\ (i\in I_z^+)},\\
 \sigma_z^+(r)&=\rAI{\alpha_z^+(r),\alpha_z^-(r)},\qquad
 \sigma_z^-(r)=\rAI{\alpha_z^-(r),\widehat\alpha_z^+(r)}.
\end{aligned}
\]
Im Term \(\widehat\alpha_z^+(r)\) liefert \(\mathrm{DN}\) aus \(z\in X_i\) die
Negation von \(z\notin X_i\). Somit stimmen sowohl die Reihenfolge
als auch die Klammerung der Prämissen mit der Aussonderungsregel
überein.
Die beiden anschließend verwendeten zusammengesetzten Quellen sind
\[
\begin{aligned}
 \kappa_z^+(r)&=\FormulaRefAuto{FranklTripleExampleDefinition}[def]
      {\FormulaRefAuto{Enthaltende Teilfamilie}[def]
      {\FormulaRefAuto{FranklTripleFilterCertificate}[thm]{\sigma_z^+(r)}}},\\
 \kappa_z^-(r)&=\FormulaRefAuto{FranklTripleExampleDefinition}[def]
      {\FormulaRefAuto{Vermeidende Teilfamilie}[def]
      {\FormulaRefAuto{FranklTripleFilterCertificate}[thm]{\sigma_z^-(r)}}}.
\end{aligned}
\]

\FormulaThmDeltaK[Explizite Gleichmächtigkeiten der Teilfamilien]{%
  \FTData\vdash
  \bigwedge_{z\in(a,b,c)}\Bigl(
     \NatSeg{\FTNum9}\EqCard\FamContaining{M_0}{z}\land
     \NatSeg{\FTNum{10}}\EqCard\FamAvoiding{M_0}{z}\Bigr)
}{FranklTripleExampleCounts}{%
  \DeltaRow{Mengen}{M_0\dsep a\dsep b\dsep c\dsep d\dsep e\dsep f\dsep g}
}
\begin{tabproofwide}
  \proofstepwidestar[1]{\FTData}{\rA}
  \proofstepwidestar[1]{\bigwedge_{(i,z)}I_{iz}}{%
    \FormulaRefAuto{FranklTripleExampleIncidence}[thm]{1}}
  \proofstepwidestar[1]{\FTDistinct{X_1,\ldots,X_{19}}}{%
    \FormulaRefAuto{FranklTripleExampleDistinct}[thm]{1}}
  \proofstepwidestar[1]{\FamContaining{M_0}{z}
       =\FTList{X_i:i\in I_z^+}}{%
    \kappa_z^+(2)\quad(z=a,b,c)}
  \proofstepwidestar[1]{\FamAvoiding{M_0}{z}
       =\FTList{X_i:i\in I_z^-}}{%
    \kappa_z^-(2)\quad(z=a,b,c)}
  \proofstepwidestar[1]{\NatSeg{\FTNum9}\EqCard\FTList{X_i:i\in I_z^+}}{%
    \FormulaRefAuto{FranklTripleListCount}[thm]
       {\rAIRepeat{36}{\operatorname{id}}{\pi_{(i,j)}(3)\ (i<j\text{ in }I_z^+)}}}
  \proofstepwidestar[1]{\NatSeg{\FTNum{10}}\EqCard\FTList{X_i:i\in I_z^-}}{%
    \FormulaRefAuto{FranklTripleListCount}[thm]
       {\rAIRepeat{45}{\operatorname{id}}{\pi_{(i,j)}(3)\ (i<j\text{ in }I_z^-)}}}
  \proofstepwidestar[1]{\NatSeg{\FTNum9}\EqCard\FamContaining{M_0}{z}
      \land\NatSeg{\FTNum{10}}\EqCard\FamAvoiding{M_0}{z}}{%
    \rAI{\rIE{\FormulaRefAuto{a=b\vdash b=a}[thm]{4},6},
          \rIE{\FormulaRefAuto{a=b\vdash b=a}[thm]{5},7}}}
  \proofstepwidestar[1]{\bigwedge_{z\in(a,b,c)}\Bigl(
     \NatSeg{\FTNum9}\EqCard\FamContaining{M_0}{z}\land
     \NatSeg{\FTNum{10}}\EqCard\FamAvoiding{M_0}{z}\Bigr)}{%
    \rAIRepeat{3}{\operatorname{id}}{8(a)\mid8(b)\mid8(c)}}
\end{tabproofwide}

\FormulaThmDeltaK[Neun gegen zehn schließt Halbhäufigkeit aus]{%
  \NatSeg{\FTNum9}\EqCard\FamContaining{M}{z}\dsep
  \NatSeg{\FTNum{10}}\EqCard\FamAvoiding{M}{z}
      \vdash\neg\HalfOcc{z}{M}
}{FranklTripleNineAgainstTen}{%
  \DeltaRow{Mengen}{M\dsep z}
  \DeltaRow{Funktionen}{F\dsep G}
}
\begin{tabproofwide}
  \proofstepwidestar[1]{\NatSeg{\FTNum9}\EqCard\FamContaining{M}{z}}{\rA}
  \proofstepwidestar[2]{\NatSeg{\FTNum{10}}\EqCard\FamAvoiding{M}{z}}{\rA}
  \proofstepwidestar[3]{\HalfOcc{z}{M}}{\rA}
  \proofstepwidestar[3]{\exists F\colon\FamAvoiding{M}{z}\inj\FamContaining{M}{z}}{%
    \rAEb{\FormulaRefAuto{Halbhäufiges Element}[def]{3}}}
  \proofstepwidestar[5]{F\colon\FamAvoiding{M}{z}\inj\FamContaining{M}{z}}{\rA}
  \proofstepwidestar[1]{\FamContaining{M}{z}\EqCard\NatSeg{\FTNum9}}{%
    \FormulaRefAuto{EqCardSymmetric}[thm]{1}}
  \proofstepwidestar[2]{\FamAvoiding{M}{z}\EqCard\NatSeg{\FTNum{10}}}{%
    \FormulaRefAuto{EqCardSymmetric}[thm]{2}}
  \proofstepwidestar[1,2,5]{\exists G\colon\NatSeg{\FTNum{10}}\inj\NatSeg{\FTNum9}}{%
    \FormulaRefAuto{EqCardInjectionTransport}[thm]{7,6,5}}
  \proofstepwidestar[9]{G\colon\NatSeg{\FTNum{10}}\inj\NatSeg{\FTNum9}}{\rA}
  \proofstepwidestar[]{\FTNum9\in\mathbb N\land\FTNum{10}\in\mathbb N}{%
    \rAI{\pi_9(\FormulaRefAuto{FranklTripleNumeralsNatural}[thm]),
       \pi_{10}(\FormulaRefAuto{FranklTripleNumeralsNatural}[thm])}}
  \proofstepwidestar[]{\FTNum9<\FTNum{10}}{%
    \FormulaRefAuto{FranklTripleListNotation}[def]
      {\FormulaRefAuto{PeanoLtAddOne}[thm]{\rAEa{10}}}}
  \proofstepwidestar[9]{\bot}{%
    \FTSource{\FormulaRefAuto{m\in\mathbb{N}\dsep n\in\mathbb{N}\dsep m < n
      \dsep F\colon\NatSeg{n}\inj\NatSeg{m}\vdash\bot}[thm]}
      {\begin{gathered}\rAEa{10},\rAEb{10},\\[-2pt]11,9\end{gathered}}}
  \proofstepwidestar[1,2,5]{\bot}{\rEE{8,9,12}}
  \proofstepwidestar[1,2,3]{\bot}{\rEE{4,5,13}}
  \proofstepwidestar[1,2]{\neg\HalfOcc{z}{M}}{\rCI{3,14}}
\end{tabproofwide}

\FormulaThmDeltaK[Keines der drei Elemente ist halbhäufig]{%
  \FTData\vdash\forall x\in\{a,b,c\}\,\neg\HalfOcc{x}{M_0}
}{FranklTripleExampleRare}{%
  \DeltaRow{Mengen}{M_0\dsep a\dsep b\dsep c\dsep d\dsep e\dsep f\dsep g\dsep x}
}
\begin{tabproofwide}
  \proofstepwidestar[1]{\FTData}{\rA}
  \proofstepwidestar[1]{\bigwedge_{z\in(a,b,c)}\Bigl(
     \NatSeg{\FTNum9}\EqCard\FamContaining{M_0}{z}\land
     \NatSeg{\FTNum{10}}\EqCard\FamAvoiding{M_0}{z}\Bigr)}{%
    \FormulaRefAuto{FranklTripleExampleCounts}[thm]{1}}
  \proofstepwidestar[1]{\neg\HalfOcc{z}{M_0}}{%
    \FormulaRefAuto{FranklTripleNineAgainstTen}[thm]
      {\rAEa{\pi_z(2)},\rAEb{\pi_z(2)}}\quad(z=a,b,c)}
  \proofstepwidestar[4]{x\in\{a,b,c\}}{\rA}
  \proofstepwidestar[4]{x\in\FTList{a,b,c}}{%
    \rIE{\FormulaRefAuto{a=b\vdash b=a}[thm]
       {\rAEb{\rAEb{\FormulaRefAuto{FranklTripleSmallLists}[thm]}}},4}}
  \proofstepwidestar[6]{x=z}{\rA\quad(z=a,b,c)}
  \proofstepwidestar[1,6]{\neg\HalfOcc{x}{M_0}}{%
    \rIE{\FormulaRefAuto{a=b\vdash b=a}[thm]{6},3(z)}}
  \proofstepwidestar[1,4]{\neg\HalfOcc{x}{M_0}}{%
    \FTCases{5;\ [6]\vdots7\ (z=a,b,c)}}
  \proofstepwidestar[1]{\forall x\in\{a,b,c\}\,\neg\HalfOcc{x}{M_0}}{%
    \rUI{\rRI{4,8}}}
\end{tabproofwide}

\section{Existenz und Rückkehr zum Hauptsatz}

Die bisherige Konstruktion setzt sieben verschiedene Mengen voraus.
Diese Voraussetzung lässt sich durch sieben konkrete natürliche Zahlen
erfüllen. Wir geben die endliche Herleitung ausdrücklich an, damit der
abschließende Satz keine bloß bedingte Aussage bleibt.

\FormulaThmDeltaK[Sieben verschiedene konkrete Zeugen]{%
  \FTDistinct{\FTNum0,\FTNum1,\FTNum2,\FTNum3,
    \FTNum4,\FTNum5,\FTNum6}
}{FranklTripleSevenWitnesses}{\DeltaRow{Mengen}{\FTNum0\dsep\cdots\dsep\FTNum6}}
\begin{tabproofwide}
  \proofstepwidestar[]{\bigwedge_{k=0}^{19}\FTNum k\in\mathbb N}{%
    \FormulaRefAuto{FranklTripleNumeralsNatural}[thm]}
  \proofstepwidestar[]{\FTNum0\neq\FTNum{j-i}}{%
    \FormulaRefAuto{FranklTripleListNotation}[def]
      {\FormulaRefAuto{PeanoZeroNeAddOne}[thm]{\pi_{j-i-1}(1)}}
        \quad(0\leq i<j\leq6)}
  \proofstepwidestar[]{\FTNum{r+1}\neq\FTNum{j-i+r+1}}{%
    \FormulaRefAuto{FranklTripleListNotation}[def]
      {\FormulaRefAuto{PeanoAddOnePreservesNe}[thm]
       {\pi_r(1),\pi_{j-i+r}(1),
          2(i,j)\ (r=0);\ 3(i,j,r-1)\ (r>0)}}
       \quad(0\leq r<i)}
  \proofstepwidestar[]{\FTNum i\neq\FTNum j}{%
    2(i,j)\ (i=0);\quad3(i,j,i-1)\ (i>0)}
  \proofstepwidestar[]{\FTDistinct{\FTNum0,\FTNum1,\FTNum2,\FTNum3,
    \FTNum4,\FTNum5,\FTNum6}}{%
    \rAIRepeat{21}{\operatorname{id}}{4(i,j)\ (0\leq i<j\leq6)}}
\end{tabproofwide}

\par\Needspace{20\baselineskip}
\hypertarget{frankl.proof.FranklRareTripleExists}{}
\noindent\textbf{Beweis des Hauptsatzes
  \FormulaRefAuto{FranklRareTripleExists}[thm].}\par\medskip
Für \((a,b,c,d,e,f,g)=(\FTNum0,\FTNum1,\FTNum2,\FTNum3,
\FTNum4,\FTNum5,\FTNum6)\) sei \(M_0\) die oben definierte Familie.
Diese Festlegung ist nur eine Substitution in bereits bewiesene Sätze.
\begin{tabproofwide}
  \proofstepwidestar[]{\FTData}{%
    \FormulaRefAuto{FranklTripleSevenWitnesses}[thm]}
  \proofstepwidestar[]{\begin{aligned}[t]
       &\Finite{M_0}\land\UnionClosedFam(M_0)\\[-2pt]
       &{}\land\{a,b,c\}\in M_0
     \end{aligned}}{%
    \FormulaRefAuto{FranklTripleExampleFamily}[thm]{1}}
  \proofstepwidestar[]{a\neq b\land a\neq c\land b\neq c}{%
    \rAIStar{\pi_{(1,2)}(1),\pi_{(1,3)}(1),\pi_{(2,3)}(1)}}
  \proofstepwidestar[]{\forall x\in\{a,b,c\}\,\neg\HalfOcc{x}{M_0}}{%
    \FormulaRefAuto{FranklTripleExampleRare}[thm]{1}}
  \proofstepwidestarbreak[]{\Finite{M_0}\land\UnionClosedFam(M_0)\land
       a\neq b\land a\neq c\land b\neq c\land
       \{a,b,c\}\in M_0\land
       \forall x\in\{a,b,c\}\,\neg\HalfOcc{x}{M_0}}{%
    \rAIStar{\rAEa{2},\rAEa{\rAEb{2}},
       \rAEa{3},\rAEa{\rAEb{3}},\rAEb{\rAEb{3}},
       \rAEb{\rAEb{2}},4}}
  \proofstepwidestarbreak[]{\exists M\,\exists a\,\exists b\,\exists c\,
     \bigl(\Finite M\land\UnionClosedFam(M)\land
       a\neq b\land a\neq c\land b\neq c\land
       \{a,b,c\}\in M\land
       \forall x\in\{a,b,c\}\,\neg\HalfOcc{x}{M}\bigr)}{%
    \rEI{\rEI{\rEI{\rEI{5}}}}}
\end{tabproofwide}
Damit liegt der vollständige Beweis in dieser Ergänzung. Im Fachband
und im Gesamtband wird nur der Hauptsatz mit seinem Verweis geführt.
